%%%%%%%%%%%%%%%%%%%%%%%%%%%%%%%%%%%%%%%%%%%%%%%%%%%%%%%%%%%%%%%%%%%%%%%%%%%%%%%%%%%%%%%%
% Projeto de Extensão da disciplina de Programação 2 / Programação 3
% Sub-projeto do ProgramAuto
% Autores:
%     Prof. Dr. Ruben Carlo Benante
%     Nome do Aluno 1
%     Nome do Aluno 2
%
% Data: 2022-02-27
%
% Assunto: escrever aqui uma breve explicação
%
%%%%%%%%%%%%%%%%%%%%%%%%%%%%%%%%%%%%%%%%%%%%%%%%%%%%%%%%%%%%%%%%%%%%%%%%%%%%%%%%%%%%%%%%


%%%%%%%%%%%%%%%%%%%%%%%%%%%%%%%%%%%%%%%%%%%%%%%%%%%%%%%%%%%%%%%%%%%%%%%%%%%%%%%%%%%%%%%%
% Para gerar o PDF use uma das 2 opções abaixo:
%
% Opção 1: com makefile, comando único.
%    $ make ext-programaX-benante-sobrenome1-sobrenome2.pdf
%
% Opção 2: com os 3 comandos diretamente:
%    $ pdflatex pext-programaX-benante-sobrenome1-sobrenome2.tex -o pext-programaX-benante-sobrenome1-sobrenome2.pdf
%    $ bibtex biblio
%    $ pdflatex pext-programaX-benante-sobrenome1-sobrenome2.tex -o pext-programaX-benante-sobrenome1-sobrenome2.pdf


%%%%%%%%%%%%%%%%%%%%%%%%%%%%%%%%%%%%%%%%%%%%%%%%%%%%%%%%%%%%%%%%%%%%%%%%%%%%%%%%%%%%%%%%
% preambulo %%%%%%%%%%%%%%%%%%%%%%%%%%%%%%%%%%%%%%%%%%%%%%%%%%%%%%%%%%%%%%%%%%%%%%%%%%%%
\documentclass[a4paper,12pt]{article} %twocolumn
\usepackage[left=2.5cm,right=2cm,top=2.5cm,bottom=2cm]{geometry}
\usepackage[utf8]{inputenc} % letras acentuadas
\usepackage[portuguese]{babel} % tradução de títulos
\usepackage[colorlinks]{hyperref}
\usepackage{algorithm} % ambiente para índice de algoritmos
\usepackage{algpseudocode} % fonte e estilo do algoritmo
\usepackage{graphicx}
\usepackage{indentfirst}
\usepackage{url}
% \usepackage{natbib}
%[noend]

\DeclareUrlCommand\email{\urlstyle{mm}} % comando para email bonito
\floatname{algorithm}{Algoritmo} % tradução da palavra algorítimo no ambiente de índice


%%%%%%%%%%%%%%%%%%%%%%%%%%%%%%%%%%%%%%%%%%%%%%%%%%%%%%%%%%%%%%%%%%%%%%%%%%%%%%%%%%%%%%%%
% capa %%%%%%%%%%%%%%%%%%%%%%%%%%%%%%%%%%%%%%%%%%%%%%%%%%%%%%%%%%%%%%%%%%%%%%%%%%%%%%%%%
\title{ProgramAuto: subtítulo}
\author{Ruben Carlo Benante \\ Autor2 \\ Autor3}

\begin{document}

\maketitle


%%%%%%%%%%%%%%%%%%%%%%%%%%%%%%%%%%%%%%%%%%%%%%%%%%%%%%%%%%%%%%%%%%%%%%%%%%%%%%%%%%%%%%%%
% resumo %%%%%%%%%%%%%%%%%%%%%%%%%%%%%%%%%%%%%%%%%%%%%%%%%%%%%%%%%%%%%%%%%%%%%%%%%%%%%%%

\begin{abstract}

\textbf{Assunto:} Ensino da Linguagem de Programação \texttt{C} / \texttt{CPP}.

% descrever em poucas palavras seu projeto aqui

\textbf{Local:} Escola Politécnica de Pernambuco - UPE/POLI

\textbf{Órgão Financiador:} N/A

\textbf{Caracterização:} Projeto de Extensão requisito da disciplina de DCExt Programação X, sub-projeto integrante do Projeto \texttt{ProgramAuto}

% Este é o fim do resumo.

\end{abstract}


%%%%%%%%%%%%%%%%%%%%%%%%%%%%%%%%%%%%%%%%%%%%%%%%%%%%%%%%%%%%%%%%%%%%%%%%%%%%%%%%%%%%%%%%
% artigo %%%%%%%%%%%%%%%%%%%%%%%%%%%%%%%%%%%%%%%%%%%%%%%%%%%%%%%%%%%%%%%%%%%%%%%%%%%%%%%
% seção de introdução %%%%%%%%%%%%%%%%%%%%%%%%%%%%%%%%%%%%%%%%%%%%%%%%%%%%%%%%%%%%%%%%%%
\section{Introdução}

% Descrever melhor seu projeto aqui

O ensino da Linguagem de Programação \texttt{C} / \texttt{CPP}, e

Diante do desafio de \ldots

O ensino de programação \ldots

Exemplo de letra em \textit{itálico é com textit}, e em \textbf{bold face é com textbf} e \texttt{tal tal tal é mono-espaçada type-writer}.


%%%%%%%%%%%%%%%%%%%%%%%%%%%%%%%%%%%%%%%%%%%%%%%%%%%%%%%%%%%%%%%%%%%%%%%%%%%%%%%%%%%%%%%%
% seção de objetivos %%%%%%%%%%%%%%%%%%%%%%%%%%%%%%%%%%%%%%%%%%%%%%%%%%%%%%%%%%%%%%%%%%%
\section{Objetivos}

\subsection{Objetivo Geral}

%Descrever o objetivo geral a ser alcançado

Criação de tal tal tal, para externalizar para a comunidade conhecimentos que tal e tal

\subsection{Objetivos Específicos}

% Listar os objetivos específicos

Como objetivos específicos este projeto visa atingir tal e tal, conforme abaixo:

\begin{itemize}
    \item Proporcionar tal e tal;
    \item Realizar tal e tal;
    \item E coisa e tal.
\end{itemize}


%%%%%%%%%%%%%%%%%%%%%%%%%%%%%%%%%%%%%%%%%%%%%%%%%%%%%%%%%%%%%%%%%%%%%%%%%%%%%%%%%%%%%%%%
% seção de justificativa %%%%%%%%%%%%%%%%%%%%%%%%%%%%%%%%%%%%%%%%%%%%%%%%%%%%%%%%%%%%%%%
\section{Justificativa}

%Justificar aqui o seu projeto.

As características de Orientação a Objetos (Encapsulamento) e RAII (Resource Acquisition Is Initialization) da linguagem C++17 são de fundamental importância para se atingir a excelência na construção de sistemas modulares, conforme denota o planejamento técnico do projeto. O encapsulamento permite a separação rigorosa entre a interface gráfica (Qt Widgets), as regras de negócio (Services) e a persistência de dados (DAOs), garantindo que a lógica de banco de dados não "vaze" para a camada visual. Além disso, a aplicação do conceito de RAII é crítica para a gestão do SQLite3, assegurando que transações e conexões com o banco sejam abertas e encerradas de forma segura, prevenindo vazamentos de recursos ou estados inconsistentes em caso de erros. \cite{Benante2008phd}.

A Figura \ref{fig:graficobarra} abaixo identifica a principal razão, exposta acima, para bla bla

\begin{figure}[ht]
\centering
\includegraphics[width=.60\linewidth]{imagens/exemplo1.png}
\caption{Exemplo de gráfico em barras}
\label{fig:graficobarra}
\end{figure}


%%%%%%%%%%%%%%%%%%%%%%%%%%%%%%%%%%%%%%%%%%%%%%%%%%%%%%%%%%%%%%%%%%%%%%%%%%%%%%%%%%%%%%%%
% seção de metodologia %%%%%%%%%%%%%%%%%%%%%%%%%%%%%%%%%%%%%%%%%%%%%%%%%%%%%%%%%%%%%%%%%
\section{Metodologia}

% Descrever como (por quais métodos) os objetivos serão alcançados.

Este projeto contará com a criação, roteiro, produção, gravação e divugação de quatro (4) vídeos tal e tal.

Os videos serão divididos nos conteúdos tal e tal.

A abordagem dada é tal e tal e coisa e tal.

%%%%%%%%%%%%%%%%%%%%%%%%%%%%%%%%%%%%%%%%%%%%%%%%%%%%%%%%%%%%%%%%%%%%%%%%%%%%%%%%%%%%%%%%
% subseção equipamentos %%%%%%%%%%%%%%%%%%%%%%%%%%%%%%%%%%%%%%%%%%%%%%%%%%%%%%%%%%%%%%%%
\subsection{Equipamentos Necessários}

%Listar os equipamentos necessários para a implementação do projeto.

Para realizar este projeto é preciso tal e tal, além de tal, e acesso à internet, e tal e tal.

Além desses equipamentos, também é preciso coisa e tal.


%%%%%%%%%%%%%%%%%%%%%%%%%%%%%%%%%%%%%%%%%%%%%%%%%%%%%%%%%%%%%%%%%%%%%%%%%%%%%%%%%%%%%%%%
% subseção com algoritmo %%%%%%%%%%%%%%%%%%%%%%%%%%%%%%%%%%%%%%%%%%%%%%%%%%%%%%%%%%%%%%%
\subsection{Implementação}

A implementação do projeto se dará de forma online, em uma \textit{playlist} numa plataforma de reprodução de videos da escolha do professor orientador, e a mesma será divulgada tanto para alunos da Escola Politécnica de Pernambuco (POLI/UPE) quanto para a comunidade em geral através dos canais de comunicação disponíveis.



%%%%%%%%%%%%%%%%%%%%%%%%%%%%%%%%%%%%%%%%%%%%%%%%%%%%%%%%%%%%%%%%%%%%%%%%%%%%%%%%%%%%%%%%
% seção Plano de Trabalho %%%%%%%%%%%%%%%%%%%%%%%%%%%%%%%%%%%%%%%%%%%%%%%%%%%%%%%%%%%%%%
\section{Plano de Trabalho}
Este projeto está subdividido em 6 etapas, que contemplam o ciclo completo de produção dos vídeos didáticos sobre simulação de sistemas computacionais, métodos estatísticos e depuração de software, e estão organizadas em ordem cronológica na Tabela \ref{tab:cronograma}.

São elas:

\begin{description}
 \item[Etapa 1: Projeto.] Organização do projeto, composição das idéias, seleção de tema, discussão com membros do grupo, divisão de tarefas e atividades administrativas relevantes.
 \item[Etapa 2: Pesquisa.] Levantamento bibliográfico e estudo dos conteúdos a serem abordados nos vídeos: método de Monte Carlo para simulação iterativa, controle robótico para simulação interativa, métodos estatísticos de comparação de resultados e testes de hipóteses com \texttt{R} ou \texttt{Octave}, e técnicas de depuração (\textit{debug}) e teste de software.
 \item[Etapa 3: Conteúdo.] Elaboração dos roteiros de cada vídeo, definindo a estrutura narrativa, a sequência lógica dos tópicos abordados e os exemplos de código a serem demonstrados em tela, garantindo linguagem didática adequada ao público-alvo.
 \item[Etapa 4: Desenvolvimento.] Gravação dos vídeos em formato \textit{screencast}, com demonstrações práticas dos códigos implementados em \texttt{C/C++}, \texttt{R} e \texttt{Octave}, seguidas da edição do material com cortes, ajustes de áudio e inserção de elementos visuais.
 \item[Etapa 5: Revisão.] Revisão do conteúdo técnico e didático dos vídeos editados, correção de eventuais imprecisões e validação do material com suporte do professor orientador antes da publicação.
 \item[Etapa 6: Validação.] Publicação dos vídeos em \textit{playlist} na plataforma definida pelo professor orientador, divulgação para os alunos da POLI/UPE e para a comunidade em geral, e entrega formal do relatório do projeto de extensão.
\end{description}
%%%%%%%%%%%%%%%%%%%%%%%%%%%%%%%%%%%%%%%%%%%%%%%%%%%%%%%%%%%%%%%%%%%%%%%%%%%%%%%%%%%%%%%%
% cronograma
%%%%%%%%%%%%%%%%%%%%%%%%%%%%%%%%%%%%%%%%%%%%%%%%%%%%%%%%%%%%%%%%%%%%%%%%%%%%%%%%%%%%%%%%
\section{Cronograma}

\begin{table}[H]
\begin{center}
 \caption{Tabela de Cronograma do Projeto de Extensão}
\begin{tabular}{|l|c|c|c|c|c|c|}
  \hline \hline
   \textbf{Etapas} & \textbf{Sem. 1} & \textbf{Sem. 2} & \textbf{Sem. 3} & \textbf{Sem. 4} & \textbf{Sem. 5} & \textbf{Sem. 6} \\ \hline \hline
   1. Projeto. & X & X &  &  &  &  \\ \hline
   2. Pesquisa. & X & X &  &  &  &  \\ \hline
   3. Conteúdo. &  & X & X &  &  &  \\ \hline
   4. Desenvolvimento. &  &  & X & X &  &  \\ \hline
   5. Revisão. &  &  &  & X & X &  \\ \hline
   6. Validação. &  &  &  &  & X & X \\ \hline \hline
\end{tabular}
\end{center}
\end{table}

%%%%%%%%%%%%%%%%%%%%%%%%%%%%%%%%%%%%%%%%%%%%%%%%%%%%%%%%%%%%%%%%%%%%%%%%%%%%%%%%%%%%%%%%
% seção de impactos alcançados %%%%%%%%%%%%%%%%%%%%%%%%%%%%%%%%%%%%%%%%%%%%%%%%%%%%%%%%%
\section{Impactos e Transferências}


%%%%%%%%%%%%%%%%%%%%%%%%%%%%%%%%%%%%%%%%%%%%%%%%%%%%%%%%%%%%%%%%%%%%%%%%%%%%%%%%%%%%%%%%
% subseção de impacto científico %%%%%%%%%%%%%%%%%%%%%%%%%%%%%%%%%%%%%%%%%%%%%%%%%%%%%%%
\subsection{Impacto Científico}

Não há impacto científico relevante (N/A).


%%%%%%%%%%%%%%%%%%%%%%%%%%%%%%%%%%%%%%%%%%%%%%%%%%%%%%%%%%%%%%%%%%%%%%%%%%%%%%%%%%%%%%%%
% subseção de impacto tecnológico %%%%%%%%%%%%%%%%%%%%%%%%%%%%%%%%%%%%%%%%%%%%%%%%%%%%%%
\subsection{Impacto Tecnológico}

Não há impacto tecnológico relevante (N/A).


%%%%%%%%%%%%%%%%%%%%%%%%%%%%%%%%%%%%%%%%%%%%%%%%%%%%%%%%%%%%%%%%%%%%%%%%%%%%%%%%%%%%%%%%
% subseção de econônimo %%%%%%%%%%%%%%%%%%%%%%%%%%%%%%%%%%%%%%%%%%%%%%%%%%%%%%%%%%%%%%%%
\subsection{Impacto Econômico}

Não há impacto econômico relevante (N/A).


%%%%%%%%%%%%%%%%%%%%%%%%%%%%%%%%%%%%%%%%%%%%%%%%%%%%%%%%%%%%%%%%%%%%%%%%%%%%%%%%%%%%%%%%
% subseção de impacto social %%%%%%%%%%%%%%%%%%%%%%%%%%%%%%%%%%%%%%%%%%%%%%%%%%%%%%%%%%%
\subsection{Impacto Social}

O projeto visa contribuir com a sociedade na forma de... bla ... bla... bla


%%%%%%%%%%%%%%%%%%%%%%%%%%%%%%%%%%%%%%%%%%%%%%%%%%%%%%%%%%%%%%%%%%%%%%%%%%%%%%%%%%%%%%%%
% subseção de impacto ambiental %%%%%%%%%%%%%%%%%%%%%%%%%%%%%%%%%%%%%%%%%%%%%%%%%%%%%%%%
\subsection{Impacto Ambiental}

Não há impacto ambiental relevante (N/A).


%%%%%%%%%%%%%%%%%%%%%%%%%%%%%%%%%%%%%%%%%%%%%%%%%%%%%%%%%%%%%%%%%%%%%%%%%%%%%%%%%%%%%%%%
% subseção de transferências %%%%%%%%%%%%%%%%%%%%%%%%%%%%%%%%%%%%%%%%%%%%%%%%%%%%%%%%%%%
\subsection{Transferências}

O ProgramAuto terá como objetivo a execução de aulas para ensinar a linguagem C aos discentes interessados. Contará com sua exposição de ensino gravada que será disponibilizada e a elaboração de relatórios a fim de cumprir com os aspectos estabelecidos no Plano de Trabalho.

Este projeto, parte integrante do ProgramAuto, visa ...

%O que mais o seu projeto agrega? O que é transferido? De onde vem? Para onde vai?


%%%%%%%%%%%%%%%%%%%%%%%%%%%%%%%%%%%%%%%%%%%%%%%%%%%%%%%%%%%%%%%%%%%%%%%%%%%%%%%%%%%%%%%%
% seção de resultados esperados %%%%%%%%%%%%%%%%%%%%%%%%%%%%%%%%%%%%%%%%%%%%%%%%%%%%%%%%
\section{Resultados Esperados}

Como resultados esperamos estabelecer, através dos vídeos, um norte para o desenvolvimento de ingressantes de disciplinas computacionais que exijam conhecimentos em programação em C++, onde muitas vezes os estudantes sentem dificuldade em encontrar com exatidão os assuntos, servindo como base inicial. No entanto, o conteúdo é livre e poderá ser acessado por todos os públicos interessados nessa área e esse é o resultado que mais esperamos como estudantes da UPE, uma instituição pública sem fins lucrativos, toda essa informação disponibilizada será gratuita.
Dado que os tópicos abordados englobam simulação interativa, interatividade, comparação de resultados, depuração, entre conteúdos que são essenciais para a área de programação, ao utilizarmos de linguagem simples, esperamos que o conteúdo seja o mais acessível o possível para sua maior disseminação e escalabilidade.
No quesito interno, esperamos, ao executar um projeto que traz um enfoque pedagógico, o aprimoramento de habilidades comunicativas, inter relacionais e criativas, focando na capacidade de explicar ideias, construindo uma linha clara de exposição de raciocínio. Além disso, como futuros engenheiros, precisamos cada vez mais aprender a desenvolver a capacidade de trabalhar em conjunto, a profissão, em inúmeros casos, exige a cooperação de várias equipes, onde entender e ser entendido é fundamental.

%%%%%%%%%%%%%%%%%%%%%%%%%%%%%%%%%%%%%%%%%%%%%%%%%%%%%%%%%%%%%%%%%%%%%%%%%%%%%%%%%%%%%%%%
% Autores %%%%%%%%%%%%%%%%%%%%%%%%%%%%%%%%%%%%%%%%%%%%%%%%%%%%%%%%%%%%%%%%%%%%%%%%%%%%%%
\section*{Detalhamento dos Autores}

%%%%%%%%%%%%%%%%%%%%%%%%%%%%%%%%%%%%%%%%%%%%%%%%%%%%%%%%%%%%%%%%%%%%%%%%%%%%%%%%%%%%%%%%
% Discentes %%%%%%%%%%%%%%%%%%%%%%%%%%%%%%%%%%%%%%%%%%%%%%%%%%%%%%%%%%%%%%%%%%%%%%%%%%%%
\subsection*{Discentes}

\begin{enumerate}
    \item \textbf{Nome Completo:} Fulano de tal
    \begin{description}
        \item [Email:] \email{blabla@poli.br}
        \item [Endereço:]
        \item [Matrícula:]
        \item [CPF:]
        \item [RG:]
        \item [Telefone:]
        \item [Currículo Lattes:] \url{http://lattes.cnpq.br/nnnnn}
    \end{description}

    \item \textbf{Nome Completo:} Fulano de tal
    \begin{description}
        \item [Email:] \email{blabla@poli.br}
        \item [Endereço:]
        \item [Matrícula:]
        \item [CPF:]
        \item [RG:]
        \item [Telefone:]
        \item [Currículo Lattes:] \url{http://lattes.cnpq.br/nnnnn}
    \end{description}

    \item \textbf{Nome Completo:} Fulano de tal
    \begin{description}
        \item [Email:] \email{blabla@poli.br}
        \item [Endereço:]
        \item [Matrícula:]
        \item [CPF:]
        \item [RG:]
        \item [Telefone:]
        \item [Currículo Lattes:] \url{http://lattes.cnpq.br/nnnnn}
    \end{description}

    \item \textbf{Nome Completo:} Fulano de tal
    \begin{description}
        \item [Email:] \email{blabla@poli.br}
        \item [Endereço:]
        \item [Matrícula:]
        \item [CPF:]
        \item [RG:]
        \item [Telefone:]
        \item [Currículo Lattes:] \url{http://lattes.cnpq.br/nnnnn}
    \end{description}

%     \item \textbf{Nome Completo:} Fulano de tal
%     \begin{description}
%         \item [Email:] \email{blabla@poli.br}
%         \item [Endereço:]
%         \item [Matrícula:]
%         \item [CPF:]
%         \item [RG:]
%         \item [Telefone:]
%         \item [Currículo Lattes:] \url{http://lattes.cnpq.br/nnnnn}
%     \end{description}
\end{enumerate}


%%%%%%%%%%%%%%%%%%%%%%%%%%%%%%%%%%%%%%%%%%%%%%%%%%%%%%%%%%%%%%%%%%%%%%%%%%%%%%%%%%%%%%%%
% Docentes %%%%%%%%%%%%%%%%%%%%%%%%%%%%%%%%%%%%%%%%%%%%%%%%%%%%%%%%%%%%%%%%%%%%%%%%%%%%%
\subsection*{Docentes}

\begin{enumerate}
    \item \textbf{Nome Completo:} Ruben Carlo Benante
    \begin{description}
        \item [Email:] \email{rcb@poli.br}
        \item [Matrícula:] 11238-0
        \item [Currículo Lattes:] \url{http://lattes.cnpq.br/3366717378277623}
    \end{description}

%     \item \textbf{Nome Completo:}
%     \begin{description}
%         \item [Email:] \email{no@no.no}
%         \item [Matrícula:]
%         \item [Currículo Lattes:]
%     \end{description}
\end{enumerate}


%%%%%%%%%%%%%%%%%%%%%%%%%%%%%%%%%%%%%%%%%%%%%%%%%%%%%%%%%%%%%%%%%%%%%%%%%%%%%%%%%%%%%%%%
% referências bibliográficas %%%%%%%%%%%%%%%%%%%%%%%%%%%%%%%%%%%%%%%%%%%%%%%%%%%%%%%%%%%
%\section*{Referências Bibliográficas}

% cite todos, mesmo os não referenciados %%%%%%%%%%%%%%%%%%%%%%%%%%%%%%%%%%%%%%%%%%%%%%%
\nocite{*}


%%%%%%%%%%%%%%%%%%%%%%%%%%%%%%%%%%%%%%%%%%%%%%%%%%%%%%%%%%%%%%%%%%%%%%%%%%%%%%%%%%%%%%%%
% se necessario %%%%%%%%%%%%%%%%%%%%%%%%%%%%%%%%%%%%%%%%%%%%%%%%%%%%%%
% troca autor and autor por autor & autor, na bibliografia. O dcu usa "and"
%\renewcommand{\harvardand}{\&} % troca and pro &. O dcu usa "and"

% Estilos de bibliografia %%%%%%%%%%%%%%%%%%%%%%%%%%%%%%%%%%%%%%%%%%%%%%%%%%%%%%%%%%%%%%
% \bibliographystyle{abnt-alf} % Estilo alfabético da ABNT. Opção [num] para estilo numérico
% \bibliographystyle{apalike}
% \bibliographystyle{dcu} %citacao como (autor and autor, ano). Parece apalike. Rev. Control. Automacao. Use com harvard
% \bibliographystyle{agsm} % padrao harvard fica (autor & autor ano).
\bibliographystyle{acm}

%%%%%%%%%%%%%%%%%%%%%%%%%%%%%%%%%%%%%%%%%%%%%%%%%%%%%%%%%%%%%%%%%%%%%%%%%%%%%%%%%%%%%%%%
% arquivo de banco de dados das referências %%%%%%%%%%%%%%%%%%%%%%%%%%%%%%%%%%%%%%%%%%%%
% renomear para o número do exercício correto
% o arquivo de bibliografia pode se chamar qualquer coisa, isso não muda o comando de gerar o PDF.
% Por exemplo para 'mybiblio.bib', use \bibliography{mybiblio} e os comandos pdflatex e bibtex continuam os mesmos identicos com exN.
\bibliography{biblio}

\end{document}
